% !TeX root = ../../Book3.tex
% 纲要标题：### 13.1 紧流形上的全局估计
% 纲要位置：几何学计划纲要.md，第 2612 行。
% 正文状态：初稿；尚未正式审核。

% 纲要内容开始
% 
% * 有限覆盖与局部估计的粘合
% * 含低阶范数的先验不等式
% * 核的有限维性
% * 核的正交补上的先验估计
% 
% 纲要内容结束

\section{紧流形上的全局估计}
\label{g3:sec:1301}

局部估计控制一个坐标片内的导数。要解整个流形上的方程，还须知道这些估计能否统一、齐次方程有多少自由度，以及去掉这些自由度后能否稳定地反解。本节在闭流形上把三个问题接起来：有限覆盖给出全局估计，紧嵌入限制核的维数，正交归一化去掉估计中最后的零阶项。

\ProseParagraphBreak

这里的闭流形指紧且无边界的光滑流形。固定光滑底度量、正体积密度、有限秩丛的度量与联络；复 Hermitian 内积仍对第一变量线性，\((u,v)_2=\int_M\langle u,v\rangle\mu_g\)。算子的阶数记为 \(d\)，流形维数记为 \(n\)。

\subsection{有限覆盖与局部估计的粘合}
\label{g3:subsec:1301:01}
先说明 Sobolev 空间与微分算子的具体对应。在有限平凡化中，光滑 \(d\) 阶算子有表达式
\[
	P=\sum_{|\alpha|\le d}A_\alpha(x)\mathrm D^\alpha,
	\qquad \mathrm D_j=-\mathrm i\partial_j.
\]
对每个整数 \(k\ge0\)，乘积求导至 \(k\) 阶时，\(Pu\) 只涉及 \(u\) 的至 \(k+d\) 阶导数及有限阶系数导数。紧集上的系数界与\cref{g3:prop:sobolev-norm-equivalence}因此给出有界映射
\[
	P:W^{k+d,p}(M;E)\longrightarrow W^{k,p}(M;F),\qquad1\le p<\infty.
\]
这个作用与分布意义的作用一致，因为光滑逼近可代入每个测试配对。椭圆性将提供反方向的估计。

\begin{theorem}[全局椭圆估计]\label{g3:thm:global-elliptic-hk}
	设 \(M^n\) 为非空闭光滑流形，\(n\ge1\)，\(E,F\to M\) 为同秩有限秩光滑实或复向量丛，固定光滑度量与联络。设 \(P:C^\infty(E)\to C^\infty(F)\) 为光滑系数的椭圆微分算子，阶数 \(d\ge1\)。对每个整数 \(k\ge0\)，存在 \(C_k\)，使
	\begin{equation}\label{g3:eq:global-elliptic-hk}
		\|u\|_{H^{k+d}}\le C_k\bigl(\|Pu\|_{H^k}+\|u\|_2\bigr)
	\end{equation}
	对所有 \(u\in H^{k+d}(E)\) 成立。若分布截面 \(u\) 满足 \(Pu\in H^k(F)\)，则 \(u\in H^{k+d}(E)\)，并满足同一估计。
\end{theorem}
\begin{proof}
	选有限平凡化 \(U_a\)，从属单位分解 \(\chi_a\)，以及 \(\psi_a\in C_c^\infty(U_a)\)，使 \(\psi_a=1\) 于 \(\operatorname{supp}\chi_a\) 附近。第十二章的分布估计\cref{g3:thm:elliptic-distribution-l2}在每个坐标片中取 \(s=k,t=0\)，得到
	\[
		\|(\chi_a u)_a\|_{H^{k+d}(\mathbb R^n)}
		\le C_a\bigl(\|(\psi_a Pu)_a\|_{H^k(\mathbb R^n)}+\|(\psi_a u)_a\|_2\bigr).
	\]
	下标表示相应丛的局部分量。局部定理已经控制了截断交换子；这一步无需把 \(P(\chi_a u)\) 错写成 \(\chi_a Pu\)。对有限个 \(a\) 求和，使用 Sobolev 范数等价及光滑乘法的有界性，即得式\eqref{g3:eq:global-elliptic-hk}。光滑稠密性使估计延伸到 \(H^{k+d}\)。

	对任意分布 \(u\)，先在每张坐标片使用同一定理的正则性部分，得到 \(u\in H^{k+d}_{\mathrm{loc}}\)。紧性与有限覆盖使其成为全局 \(H^{k+d}\) 截面，再应用已证估计。这个次序并未预先要求未知分布具有 \(L^2\) 范数。
\end{proof}

闭流形上的每个点都位于一个内部坐标片中。若 \(M\) 有边界，覆盖边界的半球不是第十二章的内部区域；边界条件和相应估计将在下一章另行建立。对变化的几何族，式\eqref{g3:eq:global-elliptic-hk}中的常数仍依赖覆盖、截断、度量、标架和所需系数界，固定流形上的有限性不自动给出族上的一致性。

\subsection{含低阶范数的先验不等式}
\label{g3:subsec:1301:02}
零阶项记录了尚未由方程控制的低频部分。要在其他可积指数中得到相同形式，先把低阶导数压到任意小倍数的高阶范数中。

\begin{lemma}\label{g3:lem:compact-sobolev-interpolation}
	设 \(M\) 为闭光滑流形，\(E\) 为光滑有限秩度量丛，\(q\ge1\) 为整数，\(1<p<\infty\)。对每个 \(\varepsilon>0\)，存在 \(C_\varepsilon\)，使
	\[
		\|u\|_{W^{q-1,p}}\le\varepsilon\|u\|_{W^{q,p}}+C_\varepsilon\|u\|_{L^p}
		\qquad(u\in W^{q,p}(E)).
	\]
\end{lemma}
\begin{proof}
	\(q=1\) 时取 \(C_\varepsilon=1\) 即可。若某个 \(q>1\)、\(\varepsilon>0\) 的结论不成立，归一化后可取 \(u_j\)，使
	\[
		\|u_j\|_{W^{q-1,p}}=1,\qquad
		\varepsilon\|u_j\|_{W^{q,p}}+j\|u_j\|_p<1.
	\]
	由\cref{g3:thm:compact-bundle-rellich}，一个子列在 \(W^{q-1,p}\) 中收敛。它的 \(L^p\) 范数趋零，故极限只能为零，却又应有 \(W^{q-1,p}\) 范数一，矛盾。
\end{proof}

\begin{theorem}\label{g3:thm:global-elliptic-wkp}
	设 \(M^n\) 为非空闭光滑流形，\(n\ge1\)，\(E,F\) 为同秩有限秩光滑度量丛，\(P:C^\infty(E)\to C^\infty(F)\) 是光滑系数、阶 \(d\ge1\) 的椭圆算子。固定整数 \(k\ge0\) 及 \(1<p<\infty\)。则
	\begin{equation}\label{g3:eq:global-elliptic-wkp}
		\|u\|_{W^{k+d,p}}\le C\bigl(\|Pu\|_{W^{k,p}}+\|u\|_p\bigr)
	\end{equation}
	对 \(u\in W^{k+d,p}(E)\) 成立；若分布截面满足 \(Pu\in W^{k,p}(F)\)，则它属于 \(W^{k+d,p}(E)\)，并满足此式。
\end{theorem}
\begin{proof}
	\(k=0\) 时，在一组内外嵌套坐标片上应用\cref{g3:thm:local-elliptic-lp-apriori}，有限相加即得结论。

	为处理高阶，先令 \(u\) 光滑。在每个稍大的坐标片中，对全部 \(|\beta|=k\) 写出
	\[
		P(\partial^\beta u)=\partial^\beta(Pu)+[P,\partial^\beta]u.
	\]
	交换子中至少一次导数落在系数上，所以其作用于 \(u\) 的阶数至多为 \(k+d-1\)。局部 \(W^{d,p}\) 估计用于 \(\partial^\beta u\)，再加上较低阶导数并以有限片拼接，得到
	\[
		\|u\|_{W^{k+d,p}}
		\le C\bigl(\|Pu\|_{W^{k,p}}+\|u\|_{W^{k+d-1,p}}\bigr).
	\]
	坐标、联络与截断带来的项也都包含在后一个范数中。上一引理取 \(q=k+d\)，选择 \(\varepsilon\) 使 \(C\varepsilon<1/2\)，移项得到所需估计。光滑稠密性处理一般 \(W^{k+d,p}\) 截面。

	分布解先由\cref{g3:thm:elliptic-distribution-lp}逐片提升为 \(W^{k+d,p}\)，再用紧性拼合并代入估计。实丛可复化后取实部，结论相同。
\end{proof}

\begin{example}\label{g3:exm:global-zero-order-term}
	连通闭流形上的非负 Laplace 算子 \(\Delta_g\) 满足
	\(\|u\|_{H^2}\le C(\|\Delta_gu\|_2+\|u\|_2)\)。常数函数的 Laplace 为零，故尚未规定均值时不能删除零阶项。相反，\(P=1+\Delta_g\) 的齐次解与自身配对给出 \(\|u\|_2^2+\|\mathrm du\|_2^2=0\)，其核为零；下面的估计将说明此时可以去掉该项。
\end{example}

\subsection{核的有限维性}
\label{g3:subsec:1301:03}
令 \(\mathcal N_P\) 表示齐次方程的光滑解空间。局部正则性使这个空间也包含全部分布齐次解，因此核的大小由全局几何决定。

\begin{proposition}\label{g3:prop:finite-elliptic-kernel}
	设 \(M^n\) 为非空闭光滑流形，\(n\ge1\)，\(E,F\) 为同秩光滑有限秩度量丛，\(P:C^\infty(E)\to C^\infty(F)\) 为光滑椭圆算子，阶 \(d\ge1\)。则
	\[
		\mathcal N_P=\{u\in\mathcal D'(E):Pu=0\}
		\subset C^\infty(E)
	\]
	为有限维空间。它是 \(L^2(E)\) 的闭子空间，所有 Sobolev 实现的核都等于它。
\end{proposition}
\begin{proof}
	分布解光滑性来自\cref{g3:cor:elliptic-smoothness}。若 \(u_j\in\mathcal N_P\) 在 \(L^2\) 中趋于 \(u\)，则对任意光滑测试截面，分部积分并取极限可得 \(Pu=0\) 分布成立；所以 \(u\in\mathcal N_P\)，核闭。

	若核无限维，逐次取与前面有限多个向量正交的单位向量，可构造 \(L^2\) 正交归一列 \(e_j\in\mathcal N_P\)。全局估计给出 \(\|e_j\|_{H^d}\le C\)。由 \(d\ge1\) 及 Rellich 紧性，一个子列在 \(L^2\) 中收敛；但不同项的 \(L^2\) 距离恒为 \(\sqrt2\)，矛盾。因此核有限维。任何 Sobolev 核都包含于分布核，而每个光滑核元素属于全部这些空间，故各实现的核相同。
\end{proof}

取 \(\mathcal N_P\) 的 \(L^2\) 正交归一基 \(e_1,\ldots,e_r\)，正交投影为
\begin{equation}\label{g3:eq:elliptic-kernel-projection}
	\Pi_Pu=\sum_{j=1}^r(u,e_j)_2e_j.
\end{equation}
\symidx{\(\mathcal N_P,\Pi_P\)：椭圆算子的光滑核及其 \(L^2\) 正交投影}
它与基的选择无关，因为 \(u-\Pi_Pu\) 与核正交。每个 \(e_j\) 光滑，Hölder 不等式还给出
\[
	\|\Pi_Pu\|_{W^{q,p}}\le C_{q,p}\|u\|_{L^p},
	\qquad q\ge0,\quad1<p<\infty.
\]
因此式\eqref{g3:eq:elliptic-kernel-projection}也定义了 \(L^p\) 上的有界有限秩投影；即使 \(p\ne2\)，我们仍用这些积分配对规定“与核正交”。

\subsection{核的正交补上的先验估计}
\label{g3:subsec:1301:04}
齐次解可以任意加到一个特解上。规定特解与核正交，既选出唯一候选，也恢复稳定估计。

\begin{theorem}\label{g3:thm:kernel-orthogonal-estimate}
	设 \(M^n\) 为非空闭光滑流形，\(n\ge1\)，\(E,F\) 为同秩光滑有限秩度量丛，\(P:C^\infty(E)\to C^\infty(F)\) 为光滑椭圆 \(d\) 阶算子，\(d\ge1\)。对整数 \(k\ge0\)、\(1<p<\infty\)，存在 \(C_{k,p}\)，使
	\begin{equation}\label{g3:eq:kernel-orthogonal-estimate}
		\|u\|_{W^{k+d,p}}\le C_{k,p}\|Pu\|_{W^{k,p}}
	\end{equation}
	对所有 \(u\in W^{k+d,p}(E)\)、\(\Pi_Pu=0\) 成立。特别地，任意 \(u\in W^{k+d,p}(E)\) 满足同一右端控制 \(\|u-\Pi_Pu\|_{W^{k+d,p}}\) 的估计。
\end{theorem}
\begin{proof}
	先证明 \(\Pi_Pu=0\) 时 \(\|u\|_p\le C\|Pu\|_p\)。若不成立，可取 \(u_j\in W^{d,p}\)，使
	\[
		\Pi_Pu_j=0,\qquad \|u_j\|_p=1,\qquad \|Pu_j\|_p\longrightarrow0.
	\]
	式\eqref{g3:eq:global-elliptic-wkp}使它们在 \(W^{d,p}\) 中有界。由 Rellich 紧性，可取子列在 \(L^p\) 中趋于 \(u\)。于是 \(\|u\|_p=1\)、\(\Pi_Pu=0\)，并且测试配对中 \(Pu=0\)。分布正则性使 \(u\in\mathcal N_P\)，所以 \(u=\Pi_Pu=0\)，矛盾。

	现在把 \(\|u\|_p\le C\|Pu\|_p\) 代入全局 \(W^{k+d,p}\) 估计，即得式\eqref{g3:eq:kernel-orthogonal-estimate}。最后 \(P\Pi_Pu=0\)，所以将同一估计用于 \(u-\Pi_Pu\) 即可。
\end{proof}

这个矛盾论证也说明常数的含义：若归一化截面能让方程残差越来越小，紧性会产生一个非零齐次解。正交条件恰好排除了这个极限。因而当一族算子的核发生跳跃时，补空间上的反解常数可能变大；指数的稳定性与这种常数的一致性是两个不同的问题。
% 纲要条目排为编号小节标题，并进入目录与书签；说明段仍保留于纲要注释中。
